% ---------------------------------------------------------------------------
% The regular sub-problem: Datadog glob <-> PromQL regex matcher.
% Numbers come from notebook 04 and data/processed/regex_glob.csv.
% ---------------------------------------------------------------------------

\section{The Regular Sub-problem}
\label{sec:regular}

Both languages filter on tag values, and both do it with a pattern language, but not the
same one. Datadog's is a glob; PromQL's is a regular expression. One is a proper subclass
of the other, so the two directions of this sub-problem are not the same question asked
twice: the main direction needs a translation that always exists, and the reverse
direction needs a decision procedure for whether one does. This section gives both, and
the decision procedure is where the course's machinery does the work.

\subsection{Globs, and what they can say}
\label{sec:globs}

\begin{definition}[Tag glob]\label{def:glob}
Fix an alphabet $\Sigma$ of tag-value characters. A \emph{glob} is a string over
$\Sigma\cup\{*\}$, written $g=s_0*s_1*\cdots*s_n$ with $n\ge 0$ stars and literal
segments $s_i\in\Sigma^{*}$. It denotes
\[
L(g)=\{\,s_0w_1s_1w_2\cdots w_ns_n \;:\; w_i\in\Sigma^{*}\,\}.
\]
Datadog documents prefix, suffix and substring wildcards and no other
construct~\citep{ddfilter}: there is no single-character wildcard, no alternation inside
a value, no character class, and no escape that would make a literal
\texttt{*}. Write $m(g)=\sum_i|s_i|$ for the total literal length.
\end{definition}

The main direction needs only the easy half, and gets it by construction.

\begin{proposition}[Globs are regular, uniformly]\label{prop:glob-regex}
Let $\mathrm{glob}(g)$ be the regular expression obtained from $g$ by replacing every
\texttt{*} with \texttt{.*} and escaping every other character. Then
$L(\mathrm{glob}(g))=L(g)$, where the regular expression is read under PromQL's whole-value
anchoring.
\end{proposition}
\begin{proof}
Induction on $n$. For $n=0$ both sides are $\{s_0\}$. For the step, \texttt{.*} denotes
$\Sigma^{*}$ and concatenation of regular expressions denotes concatenation of languages,
which is exactly how Definition~\ref{def:glob} quantifies $w_i$. Anchoring matters: it is
what makes the outermost concatenation span the whole value, and it is the fact
Proposition~\ref{prop:disjunction} rests on, read off a real server there rather than
assumed.
\end{proof}

\noindent
Every glob therefore has an exact PromQL matcher, which is why no gap code in
Section~\ref{sec:gaps} mentions globs on tag \emph{values}. The remaining direction is
the interesting one, and one structural fact decides most of it.

\begin{lemma}[Dichotomy]\label{lem:dichotomy}
For every glob $g$: if $n=0$ then $L(g)=\{s_0\}$ is a single word; and if $n\ge 1$ then
$L(g)$ contains a word of every length $\ell\ge m(g)$, and every word of $L(g)$ contains
each $s_i$ as a factor, in order.
\end{lemma}
\begin{proof}
The first case is immediate. For the second, taking $w_1\in\Sigma^{\ell-m(g)}$ and
$w_i=\varepsilon$ for $i>1$ gives a word of length $\ell$; the factor claim is
Definition~\ref{def:glob} read left to right.
\end{proof}

\noindent
So a glob language is either a singleton or unbounded in length. That is a strong
restriction, and it settles each construct the corpus refuses --- not by measurement but
by proof.

\begin{corollary}[Three things no glob says]\label{cor:no-glob}
Assume $|\Sigma|\ge 2$. Then
\begin{enumerate}
\item no glob denotes a finite language with two or more words --- in particular none
denotes $\Sigma$, so there is no single-character wildcard, and none denotes an
alternation of distinct literals;
\item no glob denotes $\Sigma^{+}$, the non-empty values;
\item for $\emptyset\neq A\subsetneq\Sigma$, no glob denotes $A^{+}$ --- a star cannot be
restricted to part of the alphabet.
\end{enumerate}
\end{corollary}
\begin{proof}
(1) A finite language with two words is neither a singleton nor unbounded, so
Lemma~\ref{lem:dichotomy} excludes it. (2) $\varepsilon\notin\Sigma^{+}$ forces
$m(g)\ge 1$, since $m(g)=0$ makes $g$ a string of stars and $\varepsilon\in L(g)$; and
$\Sigma^{+}$ contains words of length $1$, so $m(g)\le 1$. Hence $m(g)=1$: $g$ has a
single literal character $c$ and at least one star. By Lemma~\ref{lem:dichotomy} every
word of $L(g)$ contains $c$, while $\Sigma^{+}$ contains the one-character words of
$\Sigma\setminus\{c\}$, which is non-empty. (3) The same argument gives $m(g)=1$ with one
literal $c$, which must lie in $A$ since every word of $L(g)$ contains it. A single
literal and at least one star leaves three shapes: $c{*}$, ${*}c$ and ${*}c{*}$. In the
first and third $L(g)$ contains $c\sigma$ for every $\sigma\in\Sigma$; in the second it
contains $\sigma c$ for every $\sigma$. Either way it contains a word using a
$\sigma\notin A$, which $A^{+}$ does not, and $\Sigma\setminus A\neq\emptyset$.

Taking $c\sigma$ in every case, as an earlier version of this proof did, is wrong for
$g={*}c$: that language has no word beginning with $c$ unless $c$ repeats.
\end{proof}

\subsection{Deciding whether a regex has a glob}
\label{sec:decide}

Corollary~\ref{cor:no-glob} rules particular shapes out. The general question --- given a
regular expression, is there \emph{any} glob denoting the same language? --- is settled by
a single observation about shortest words.

\begin{proposition}[Glob-expressibility is decidable]\label{prop:decidable}
Let $r$ be a regular expression with $L(r)\neq\emptyset$, and read \texttt{*} as matching
the empty string. It is decidable whether some glob $g$ has $L(g)=L(r)$, and a witness is
produced when one exists.
\end{proposition}
\begin{proof}
Suppose $L(r)=L(g)$ for $g=s_0*\cdots*s_n$. Taking every $w_i=\varepsilon$ gives the word
$s_0s_1\cdots s_n\in L(g)$ of length $m(g)$, and every other word of $L(g)$ has some
$w_i\neq\varepsilon$ and so is strictly longer. Hence $s_0s_1\cdots s_n$ is the shortest
word of $L(g)$ \emph{and the only word of that length}.

So the procedure is: construct a DFA for $r$ --- Rich's regular-expression-to-FSM
construction (Thm 6.1) followed by the subset construction (Thm 5.3) --- and find its
shortest accepted
word $w$ by breadth-first search from the start state, together with the number of
accepting paths of that length. If $L(r)$ has no unique shortest word, no glob exists. If
it does, then $w=s_0s_1\cdots s_n$, so $g$ is $w$ with stars inserted at some subset of
its $|w|+1$ gaps: at most $2^{|w|+1}$ candidates. Each is settled by deciding
$L(r)=L(\mathrm{glob}(g))$, which is decidable by Rich Thm 9.5. The search is finite and
exhaustive, so it answers either way.
\end{proof}

\noindent
Two remarks on making this run. Breadth-first search is not an optimisation but the
method: enumerating candidate words in order instead, as a general-purpose regular-language
library does, does not terminate in practice on a corpus regex whose shortest word is
sixty characters over a twenty-character alphabet. And the $2^{|w|+1}$ candidates can be
pruned without losing completeness, because inserting stars only grows the language: any
star belonging to a solution must, on its own, keep the candidate inside $L(r)$, so gaps
failing that single subset test cannot appear in any solution at all. On the corpus this
leaves a handful of candidates rather than $2^{61}$.

\subsection{Unions, and the same fact in both directions}
\label{sec:glob-unions}

A single glob is the wrong target. A Datadog scope can disjoin over one tag key, as
$k \in (g_1,\dots,g_B)$ or as $k{:}g_1 \vee \cdots \vee k{:}g_B$, and
Proposition~\ref{prop:disjunction} showed that such a scope is exactly one PromQL
alternation matcher. Read in this direction the same proposition says that the reverse
translation of a regular expression should be sought as a \emph{finite union} of globs,
and $\mathtt{LIST|GET}\mapsto k \in (\texttt{LIST},\texttt{GET})$ is then immediate. The
two directions are not two results; they are one fact used twice.

The union must be bounded, and the bound is part of the claim rather than an
implementation detail. Every finite language is a union of singleton globs, one per word,
so boundedness is what stops the question from being trivial --- and it is not only a
matter of taste. By Corollary~\ref{cor:no-glob} no glob denotes $\Sigma^{+}$, yet
$\bigcup_{c\in\Sigma}L(c*)=\Sigma^{+}$, so the non-empty values \emph{are} a finite union
of globs: one for every character of the alphabet. A scope is a piece of text somebody
maintains, so a union of a million branches is not a translation. We fix $B=64$ and
report anything needing more as a refusal, naming the bound.

Whether glob-union-expressibility is decidable without a bound we leave open. The shortest-word
argument does not transfer: each glob of a minimal union contributes its own shortest
word, and nothing bounds those by the shortest word of the union, so the candidate space
is no longer finite. Within the bound the procedure used here is sound but not known to be
complete --- it proposes a union from the regular expression's syntax and then decides
equivalence, so what it returns is exact, while a refusal is in general a statement about
the proposer. Proposition~\ref{prop:decidable} removes that caveat for the single-glob
case, and Section~\ref{sec:glob-corpus} uses it to check the refusals rather than trust
them.

\subsection{What the corpus contains}
\label{sec:glob-corpus}

\begin{table}[htbp]\centering\footnotesize
\begin{tabular}{@{}lrrrr@{}}
\toprule
Outcome & Distinct & Share & Occurrences & Share \\
\midrule
One glob & 61 & 33.2\% & 309 & 29.7\% \\
A finite union of globs & 73 & 39.7\% & 344 & 33.0\% \\
Neither ($\bot$) & 50 & 27.2\% & 388 & 37.3\% \\
\emph{Template-valued, not a regex} & 174 & --- & 7,546 & --- \\
\bottomrule
\end{tabular}

\caption{Every \texttt{=\textasciitilde{}} and \texttt{!\textasciitilde{}} matcher in the
PromQL half of the corpus. Shares are of the concrete regular expressions; the
template-valued row is counted but not classified, because the pattern is not in the query
text.}
\label{tab:regex-glob}
\end{table}

The first result is not about regular languages. Of $8{,}587$ regex matchers,
$7{,}546$ --- $87.9\%$ --- have a Grafana template variable as their pattern, as in
\texttt{\{pod=\textasciitilde{}"\$pod"\}}. The value is chosen when somebody looks at the
dashboard, so there is no pattern to convert; the normalisation of
Section~\ref{sec:corpus} deliberately leaves these alone, because the matcher already
parses with the variable in place. For these the reverse direction is not hard but
\emph{statically undefined}, and that is a property of how dashboards are written rather
than of either language.

Among the $1{,}041$ occurrences whose pattern is concrete, Table~\ref{tab:regex-glob}
gives the three-way split and Table~\ref{tab:regex-glob-examples} what each outcome
looks like. A glob exists for $653$ of them ($62.7\%$), but for more than
half of those only as a union: one glob covers $309$ occurrences ($29.7\%$) and a bounded
union covers $344$ ($33.0\%$). Those counts rose in the seventh reading of the apparatus,
and the caption of Table~\ref{tab:regex-glob} is why: it says ``every \texttt{=\textasciitilde{}}
and \texttt{!\textasciitilde{}} matcher'', and the scanner behind it accepted only
double-quoted patterns. PromQL writes a string three ways, its own grammar in
Section~\ref{sec:grammars} says so, and $114$ matcher occurrences sat in the other two. Treating a single glob as the target --- the obvious
reading of ``does this regex have a glob'' --- would therefore have understated the
reverse direction by $33.0$ points of occurrence coverage, and the construction that
recovers them is the one Section~\ref{sec:translation} already proved for the other
direction.

\begin{table}[htbp]\centering\footnotesize
\begin{tabular}{@{}lrrl@{}}
\toprule
Construct with no glob & Distinct & Occurrences & Shortest example \\
\midrule
No non-empty wildcard (\texttt{.+}) & 21 & 238 & \texttt{.+} \\
No single-character wildcard (\texttt{.}, \texttt{5..}) & 26 & 140 & \texttt{5..} \\
No bounded alphabet under a star (\texttt{[a-z]+}) & 3 & 10 & \texttt{5[0-9]+} \\
\bottomrule
\end{tabular}

\caption{The constructs behind the $388$ refused occurrences. Each is ruled out as a
single glob by Corollary~\ref{cor:no-glob}; only the third is also beyond a bounded
union of them.}
\label{tab:regex-glob-reasons}
\end{table}

\begin{table}[htbp]\centering\footnotesize
\begin{tabular}{@{}llr@{}}
\toprule
PromQL regex & Datadog scope value & Occurrences \\
\midrule
\texttt{\textasciicircum{}k8s\_.*} & \texttt{k8s\_*} & 60 \\
\texttt{.*external-dns.*} & \texttt{*external-dns*} & 36 \\
\texttt{.*pod.*} & \texttt{*pod*} & 24 \\
\texttt{LIST|GET} & \texttt{GET, LIST} & 41 \\
\texttt{ext.*|xfs} & \texttt{ext*, xfs} & 34 \\
\texttt{proxy|attach|log|exec|portforward} & \texttt{attach, exec, log, portforward, p\dots{}} & 32 \\
\texttt{.+} & $\bot$ No non-empty wildcard & 133 \\
\texttt{(/dev/)?(mmcblk.p.+|nvme.+|rbd.+|\dots{}} & $\bot$ No non-empty wildcard & 49 \\
\texttt{5..} & $\bot$ No single-character wildcard & 40 \\
\bottomrule
\end{tabular}

\caption{The three most frequent regular expressions of each outcome.}
\label{tab:regex-glob-examples}
\end{table}

Table~\ref{tab:regex-glob-reasons} is the part this section can be held to, and it has
to be read against the target of \S\ref{sec:glob-unions} rather than against a single
glob. Three constructs account for every refusal, and each is an instance of
Corollary~\ref{cor:no-glob} \emph{or} of one of the two facts below: a single-character
wildcard such as \texttt{5..} would need a
glob for $\Sigma$, which part~(1) forbids; \texttt{.+} would need one for $\Sigma^{+}$,
forbidden by part~(2); and \texttt{5[0-9]+} would need a star confined to the digits,
forbidden by part~(3).

``Each is an instance'' is the part that needed care, and an earlier version of this
sentence overstated it. The corollary is about three \emph{specific} languages and is
closed under neither concatenation with a literal nor union, while most refused corpus
patterns are of those shapes: \texttt{dm-.+} is not $\Sigma^{+}$, and
\texttt{5[0-9]+} is not $A^{+}$, because $A^{+}$ contains words of length one and that
language's shortest word has three characters. Two facts close the gap, and both are
already in this section rather than new.

The first is the uniqueness step inside Proposition~\ref{prop:decidable}'s proof: if
$L(r)=L(g)$ then $L(r)$ has a \emph{unique} shortest word. Most of the patterns the
corollary does not reach fail that immediately --- \texttt{dm-.+} has shortest words
\texttt{dm-} and nothing else of that length only if the suffix can be empty, which
\texttt{.+} forbids, so its shortest words are \texttt{dm-}$\sigma$ for every $\sigma$,
and there are $|\Sigma|$ of them. The second is the generalisation part~(3) is an
instance of: \emph{if $L$ is infinite and some $\sigma\in\Sigma$ occurs in no word of
$L$, then $L$ is no union of finitely many globs}, because an infinite glob language
needs a star and a star admits every character. That is what the prose below has always
said, and stating it as a fact rather than as an aside is what the $10$ unconditional
occurrences actually rest on.

The three are not equally final, and saying so costs less than leaving it implied. The
target here is a union of at most $B$ globs, and the first two constructs \emph{are}
finite unions: $\Sigma^{+}$ is $\bigcup_{c\in\Sigma}L(c*)$ and $\Sigma$ is the union of
its one-character words, so each needs one branch per alphabet character, or
$|\Sigma|^{2}$ of them for \texttt{5..}. Those refusals therefore rest on the bound as
well as on the corollary, and on an alphabet larger than $B$, which this paper does not
pin down --- the proposer refuses them structurally rather than by counting branches.
Only part~(3) is unconditional: an infinite language needs a star, a glob's star admits
every character, and no union of finitely many such globs confines it to the digits. So
$378$ of the $388$ occurrences are refused by a theorem about single globs plus a bound
this paper chose, and $10$ by a theorem about globs outright. For
the single-glob question we can also say it the other way round, by measurement rather
than by proof: re-deciding every concrete regular expression with the complete procedure
of Proposition~\ref{prop:decidable} reaches \emph{all} $184$ distinct patterns --- the
counts in this paragraph are of patterns rather than of the occurrences above --- and
disagrees with the syntax-directed proposer on \emph{none} of them. It recovers the
same glob for all $61$ of the single-glob cases, and finds no glob for all $50$ refusals
and all $73$ unions. Nothing is left undecided: no pattern exceeds the search's own
candidate bound, so the agreement is over the whole population rather than over a
reachable part of it.

An earlier version of this paragraph reported $168$ of $183$ and recorded the remainder
as not reached, which was a statement about a clock rather than about the procedure. The
script takes a time budget, that run was given $2{,}400$\,s, and four patterns of one
family --- a parenthesised alternation inside a long literal, as in
\texttt{cnpg\_pg\_stat\_(bgwriter|checkpointer)\_checkpoints\_timed} --- cost $91$ to
$176$ seconds each when timed on an idle machine, so the budget expired partway through
them. Run without one the cross-check takes $4{,}725$\,s and finishes. What makes those
patterns expensive is worth naming, because it is the procedure working rather than
struggling: the shortest word of such a pattern is forty-odd characters, so the pruned
space of star insertions is large, and the search has to \emph{exhaust} it to conclude
that no single glob exists. The costly cases are the negative ones, and they are exactly
the cases a structural proposer's refusal would otherwise have to be taken on trust. The run is cached
(\texttt{scripts\slash crosscheck\_globs.py}), and \texttt{tests\slash test\_regular.py}
fails if a disagreement ever appears, if any pattern goes unreached, or if the two sides
stop sharing a population.

One qualification on ``complete''. The search is complete for the regular expressions its
parser accepts, which is the dialect of the equivalence library rather than all of
PromQL's: a pattern carrying an interior anchor, or one whose escapes the structural
proposer normalises first, is refused by the search where the proposer answers. No corpus
pattern is of either shape --- all $184$ were reached --- so the agreement reported above
is over the whole population, but the word means complete relative to that dialect.

\begin{remark}[The largest refusal turned on an undocumented detail]%
\label{rem:star-empty}
$238$ of the $388$ refused occurrences are \texttt{.+} or a pattern containing it,
including the corpus's most common concrete regular expression (\texttt{.+}, $133$
occurrences) and its most elaborate one
(\texttt{(/dev/)?(mmcblk.p.+|nvme.+|\dots{})}, $49$). They are refused by
Corollary~\ref{cor:no-glob}(2) --- but that part of the corollary assumed \texttt{*}
matches the empty string, and Datadog does not document whether it does. The two readings
cannot both be wrong and cannot both be right: a glob has exactly one wildcard, which
either requires a character or does not, so under one reading \texttt{.+} has no glob and
under the other \texttt{.*} has none. The corpus uses \texttt{.+} $133$ times and
\texttt{.*} $9$ times, so the reading is worth $124$ occurrences either way.

It is settled, and in favour of the reading taken here. Submitting two series of one
metric whose \texttt{probe} tag is \texttt{web-} and \texttt{web-1}, and filtering on
\texttt{probe:web-*}, returns both: the wildcard matches the empty string. So
\texttt{.*} has a glob, \texttt{.+} has none, and the figures above stand as computed
(\texttt{scripts\slash trial\_checks.py e5}, cached and regression-tested).
\texttt{regular.py} still records the reading in a single constant, so the measurement
can be redone under the other one if Datadog's behaviour changes.
\end{remark}

\noindent
The practical conclusion for the reverse direction is narrower than
Table~\ref{tab:regex-glob} suggests on its own. A Datadog scope cannot say what
\texttt{5..} says, and cannot say what \texttt{.+} says under the reading taken here, in
any scope a person would maintain; both are ordinary things to write in PromQL. That is
the kind of gap this paper set out to distinguish from the other kind, and in this
corner of the problem no refusal is unfinished work: each is either a theorem about
globs or that theorem together with the branch bound, and neither is removed by
building more.
